% TOP_PROBLEM: {"problemId": "problem.grothendieckkatz-p-curvature-conjecture", "canonicalTitle": "Grothendieck–Katz p-Curvature Conjecture", "releaseRank": 118, "primaryDomain": "number_theory_arithmetic_geometry", "primaryDomainLabel": "Number theory & arithmetic geometry", "displayStatus": "open_with_solved_subcases", "releaseStatus": "open_with_solved_subcases"}
% !TeX program = lualatex
% Definition and sources only; this file is not an LLM solution attempt.
\documentclass[11pt]{article}
\usepackage[margin=25mm]{geometry}
\usepackage{fontspec}
\setmainfont{Latin Modern Roman}
\usepackage{amsmath,amssymb,mathtools}
\usepackage{unicode-math}
\setmathfont{Latin Modern Math}
\usepackage{xurl,hyperref}
\hypersetup{colorlinks=true,urlcolor=blue}
\setcounter{secnumdepth}{0}
\setlength{\parindent}{0pt}
\setlength{\parskip}{5pt}
\setlength{\emergencystretch}{3em}
\begin{document}
\section*{\texorpdfstring{Grothendieck–Katz p-Curvature Conjecture}{Grothendieck–Katz p-Curvature Conjecture}}\label{118}
\subsection{Definitions and mathematical statement}
For a connection in characteristic $p$, its $p$-curvature is the map on derivations
\[
 \psi_p(D)=\nabla_D^{\,p}-\nabla_{D^{[p]}},
\]
where $D^{[p]}$ is the $p$-fold derivation. Start with an algebraic integrable connection in characteristic zero and a model over a finitely generated ${\mathbb{Z}}$-algebra. If the $p$-curvature vanishes for almost all reductions in the source's sense, the conjecture predicts finite monodromy. Equivalently, after a finite étale cover the bundle with connection has a full horizontal basis. Vanishing at finitely many primes does not suffice, and a choice of model is part of expressing the arithmetic hypothesis.
\par\smallskip\textit{Formulation and concept references:} \href{https://www.math.uchicago.edu/~farb/papers/pcurv.pdf}{[S1]}.

\subsection{Short English statement}
Does vanishing $p$-curvature at almost every prime force a characteristic-zero differential equation to have only algebraic solutions and finite monodromy?
\subsection{Sources}
\begin{itemize}
\item[S1]Benson Farb and Mark Kisin, The Grothendieck–Katz Conjecture for Certain Locally Symmetric Varieties.
\url{https://www.math.uchicago.edu/~farb/papers/pcurv.pdf}.
\textit{Formulation source.}
\end{itemize}
\end{document}
